\section{Related Work}

\subsection{Image manipulation detection and localization}
Image-forensics methods exploit inconsistencies in acquisition, compression, noise, and editing traces, while learned approaches estimate manipulation scores or suspicious regions directly from image content \cite{farid2009imageforgery,verdoliva2020mediaforensics}. Research has addressed splice detection through learned self-consistency and rich manipulation features, and later systems have combined anomaly cues with deep localization architectures \cite{huh2018fighting,zhou2018rich,wu2019mantra,yu2024diffforensics}. MVSS-Net uses noise-sensitive and boundary-artifact views with multi-scale supervision, whereas TruFor fuses high- and low-level clues and outputs image-level integrity scores and localization maps \cite{chen2021mvss,guillaro2023trufor}. These systems motivate our choice of two frozen detectors with different designs. Our study does not retrain either method or propose a competing architecture; it measures how their released implementations behave under a common evaluation protocol.

\subsection{Benchmark data and evaluation scope}
Public datasets provide complementary but non-equivalent evaluation conditions. Columbia Uncompressed isolates splicing in uncompressed images and includes a dataset-specific reference annotation; CASIA provides authentic and manipulated examples spanning splicing and copy-move; COVERAGE was designed for copy-move forgery evaluation \cite{hsu2006columbia,dong2013casia,wen2016coverage}. Their source composition, image processing and annotation semantics differ. Accordingly, aggregate pixel scores should be interpreted with the provenance and validity of each reference mask in view. Generalization-focused synthetic-image detectors and broader benchmarks such as GenImage and DF40 study a related but distinct image-authentication problem, while ForensicHub addresses a wider set of image-forensics tasks \cite{ojha2023universal,zhu2023genimage,yan2024df40,du2025forensichub}. The present study is restricted to local image manipulation and the three named datasets.

\subsection{Explainability, audit records and reproducibility}
General attribution techniques such as LIME, Integrated Gradients, SHAP and Grad-CAM explain model outputs through feature relevance or saliency estimates \cite{ribeiro2016lime,sundararajan2017axiomatic,lundberg2017unified,selvaraju2017gradcam}. Recent systems extend image-forgery detection with multimodal explanations, textual rationales and localization, including FakeShield, SIDA and X2-DFD \cite{xu2024fakeshield,huang2025sida,chen2025x2dfd}. Such work makes explanation outputs available for inspection, while reproducibility of a prediction is a separate property: it requires identification of the input, model artifact, execution settings and derived output, together with a procedure for replay.

Digital-evidence guidance emphasizes documenting methods, findings and limitations, and forensic evaluation initiatives consider performance in operationally relevant settings \cite{swgde2025imageauth,nist2025guardians,nistopenmfc}. Broader discussions of AI in digital forensics and validation of forensic evaluation methods similarly motivate explicit scope and error characterization \cite{casey2022digitalforensics,ramos2017lrvalidation}. ForenX-Bench operationalizes a limited technical audit around these concerns: it records the evaluated artifacts and conditions, verifies source hashes, and tests deterministic replay in an independent process. It does not test examiner performance, validate a legal opinion, or equate reproducible output with causal evidence.
